% Supplementary details for the security analysis.
\section{Details of the Original-Execution Connection}
\label{supp:original-execution}

\noindent\textbf{Earliest conflicting polka.}
Fix a height and a commit certificate for a nonempty value $v$ in round $r$, and select two honest signers $K$ of that certificate. Each locks $v$ before it issues its precommit. The certificate can be assembled after these signers have advanced, so its observation time is irrelevant.

Suppose a higher-round polka for a different value exists, counting nil if it permits unlocking. Select the smallest round $s>r$ with such a polka and, within $s$, the earliest event $\tau$ at which its distinct prevotes exist. Its three voters include a member $m\in K$. Honest rounds increase, so the round-$s$ prevote of $m$ follows its round-$r$ precommit, and without an unlocking justification $m$ must prevote its locked value. Relocking on $v$ in intervening rounds preserves this restriction. A conflicting justification from a round below $s$ contradicts the choice of $s$, and one from round $s$ must have existed before $m$ cast the vote counted at $\tau$, which contradicts the choice of $\tau$. Evidence from above $s$ cannot precede a later round-$s$ vote. No such polka exists, and a higher-round conflicting commit would require one. Same-round conflicting commits require an honest double precommit. The argument rests on the stated lock-transition and round-monotonicity rules and does not need timely receipt of commit evidence.

\noindent\textbf{Pairing, validation, and execution.}
Induct on candidate updates. Local construction pairs a body with its original parts, and remote completion rebuilds the body from checked parts. Moving Locked or Valid transfers both, and switching targets clears the old body. Historical installation leaves these active candidates alone. A nonempty precommit follows a matching polka. An honest prevote has a justification in actual proposal validation or an earlier lock, and tracing the latter lowers the lock round until it reaches validation. The final full-identity guard fixes the original transactions, height, time, and other application context before Finalize. Synchronization and replay supply the same original material, so same-height retries cannot introduce a second original body through the stable-hash cache within this call boundary.

\noindent\textbf{Write sets.}
The proofs of the main text rely on the following write sets. A compensation decision writes the issuer reserve account, the obligation status, the Promise and coverage record (Paid and Remaining), grant usage (Spent), the public gap, the consumer's payment record (Pending and fee stages), reward and refund balances and fee outputs, the removal of the due-index entry, and the decision record with its decision-index and pending-representation entries. The execution of a source after compensation writes the repayment: the issuer reserve account, the obligation status (Recovered), the coverage record (Recovered), and grant usage (Spent). It leaves the decision record unchanged. A resubmission job writes the member's local outbox only. A representation command or batch writes the logical revision, the head command, the task or batch task with its per-output references, the queue entry, and the removal of pending-representation entries. Installation writes local block-store versions and its own progress cursor only. No representation or installation write belongs to the economic projection.

\noindent\textbf{Convergence for a fixed target set.}
Fix an original body and a set of decided slots. The replacement message of each slot is a function of the network and the output, so any packaging replaces the same slots by the same messages. If the chameleon relation has a unique valid opening for each fixed message and commitment, as holds for correctly generated parameters with a public exponent coprime to the group order, each slot opening and each part opening is the unique solution of its relation. Part contexts depend on the height and the part index and are independent of the revision counter. Any grouping or order of representation commands that covers the set therefore ends in the same body and part openings. Single-item commands, one batch, and batches of intermediate sizes differ in command identities, revision numbers, and application hashes. The statement concerns final bytes and depends on parameter generation and encoding as stated.

\noindent\textbf{Follower updates.}
For a follower that starts from the same local state, configuration, and cursor, equal fixed payment fields and authenticated results induce equal economic updates whenever both representations decode. The mutable funding reference and opening select no Paid, Pending, or fee effect. Decisions are authenticated in full, and the execution-result commitments authenticate their effects and fee outputs. Economic observation thus stays separate from the independent verification of the latest mutable historical body.

\noindent\textbf{Relations used in the analysis.}
Three relations compare different objects. Representation inertness compares the economic projection before and after one representation or installation step, and it projects away command identities, revision counters, and materialization tasks. Fixed-target convergence compares the final body and part openings that different groupings and orders of representation commands produce for one original body and one decided set. Replay equality fixes the entire original command sequence and its block contexts, and deterministic execution then reproduces each application hash. Cross-history comparisons fix the economic commands, public times, and historical target coordinates, so projection equality does not identify entire application states and does not permit deleting or renumbering blocks. None of the three relations compares histories that order parent arrivals and compensations differently; the delay-neutrality theorem of the main text makes that comparison under its own hypotheses and equates final CAL holders and reserve balances, not application hashes or revision metadata.

\noindent\textbf{Local invalidation write set and replay.}
Invalidation is evaluated only while applying an authenticated block
entry whose execution result is successful and applied; \texttt{INSTALL}
does not trigger it.  Before the follower overwrites a spend record, it
reads the previous candidate and verifies the immutable approval and
the exact principal (final or certified) or user-fee input instance.  Existing observation,
installed certificate material, an outbox entry, settlement, pending
obligations, compensation, recovery, or a fee effect is treated as an
accounting contradiction rather than as releasable state.

The additional invalidation write set contains only the affected local resource slices
and the approval's local progress.  For each debit it subtracts
$\mathsf{Cap}-\mathsf{Applied}$ from Reserved, adds the same amount to
Available, sets the corresponding cumulative Applied entry to Cap,
and records Invalidated, the first triggering public fact, and its
height.  The immutable approval and unrelated input locks remain.
Overlay visibility makes later conflicting inputs or transactions in
the same block no-ops for restoration; the follower cursor suppresses
block replay, and the cumulative Applied vector prevents a second
credit after restart.  The member-store schema is advanced so that an
older binary cannot silently ignore the invalidated terminal state. These writes commit with the successful consumption record and follower cursor. Schema 8 rejects earlier member stores; experiments use fresh directories rather than a migration.

A fully certified source that is permanently rejected by the global
Intent binding while spending disjoint inputs is outside this
invalidation rule: its QC-backed responsibility, input locks, and
signing reservation remain.  The current relay may continue retrying
such a source; terminating those retries is a separate engineering
follow-up and cannot release the certified liability.

\subsection{Executable Interface and Fee Evidence}
\label{supp:interface-fee-evidence}
The artifact's 2026-10-04 interface audit records source fingerprints,
commands, and assertion logs. It adds tests to checkout \texttt{863952d};
node production logic is unchanged from the measured \texttt{f856c7b}
baseline. Local regressions use Go 1.27.1 on Windows amd64 and supply
correctness evidence, not replacement performance measurements.

\noindent\textbf{Identity-sensitive transitions.}
The pinned CometBFT overlay replaces eight body-hash comparisons by
\texttt{matches\allowbreak BlockID}, which checks the body's hash and its paired part-set
header. The remaining Finalize hash check follows a part-header check;
the vote-extension caller is already gated by the full-identity match.
\texttt{Test\allowbreak Stable\allowbreak Hash\allowbreak Identity} exercises actual consensus state methods for
commit-target replacement, stale parts, exact finalization, Locked/Valid
pairing, POL rounds, and a quorum preceding or following the available body.
The polka-before-parts case also starts at Propose: completion causes one
precommit for the exact identity, and subsequent attempts cause neither a
late prevote nor another precommit. The separate nil-precommit case checks
that late parts cannot produce a second vote. This tests the added
current-round-polka transition as well as the comparison guards.

\noindent\textbf{Original material and application authentication.}
\texttt{Test\allowbreak Consensus\allowbreak Rejects\allowbreak Relabelled\allowbreak Opening} passes a compatible adaptation
relabelled as revision zero through the consensus message's
\texttt{Validate\allowbreak Basic} and \texttt{Msg\allowbreak From\allowbreak Proto}: original admission rejects
the nonoriginal opening. Part-set inclusion verification alone accepts
the relabelled part under the signed root; original-opening rejection
occurs at the consensus message boundary. The real BlockStore rewrite/replay and catch-up
tests retain and serve original bodies and parts after revision.
\texttt{Test\allowbreak Public\allowbreak Authorization\allowbreak At\allowbreak ABCIEntry} checks rejection by
\texttt{CheckTx}, \texttt{ProcessProposal}, and \texttt{FinalizeBlock}, with
no business-state change, after warming a valid verification cache. In its
forged-QC case, decoding, owner authorization, admission preparation, and
summary binding pass before QC authentication returns \texttt{ErrAuth}; the
unrelated-QC and forged-owner cases reject at earlier guards. Separate cache
tests require identical bytes for reuse and recheck current ledger validity.

\noindent\textbf{Escrow and closure.}
\texttt{PrepareDirectVector} checks the base fees plus one repair fee per
certified input with checked arithmetic. \texttt{feeStage} applies cumulative
escrow deltas, \texttt{endObligation} permits one repair charge on leaving
Open, and \texttt{closeFee} accounts for actual paid fees within the admitted
cap. \texttt{Test\allowbreak Fee\allowbreak Closure\allowbreak Multiple\allowbreak Obligations} has twelve rule-level cases:
two fee modes, zero to two compensated inputs, and minimum or surplus escrow.
They assert escrow and asset conservation after each transition, exact final
charges and refunds, closure at zero Pending, and idempotent duplicates.
Two further cases in \texttt{Test\allowbreak Fee\allowbreak Closure\allowbreak Fulfillment\allowbreak Before\allowbreak Compensation}
use distinct parents: one arrives normally before compensation of the other
closes the consumer, and the late second source leaves its closed escrow
unchanged. Public-decision and multi-slot tests separately cover the
application boundary. The member-following suite compares local escrow with
the public payment state after compensation and source arrival, in addition
to checking cumulative restoration and non-recycling of spent FUEL.

\noindent\textbf{What the observations establish.}
These tests connect the proof's transition assumptions to implementation
entry points; they do not replace the locking or accounting arguments.
The no-obligation decision guard rejects without changes or effects, matching
the service contract's absence of certificate-only redemption. The final
mixed-load capital audit reads the recorded counters and resource endpoints
of eight members per run: all three runs have zero resource-insufficiency
counters and zero final CAL/work reservations, while actual FUEL charges
remain retained. This endpoint audit is separate from the analytic coverage
of unpublished certificates and from the workload conditions for ongoing
admission.
